\documentclass[11pt]{article}
\usepackage[T1]{fontenc}
\usepackage[margin=1in]{geometry}
\usepackage{amsmath,amssymb,amsthm}
\usepackage[hidelinks]{hyperref}

\newtheorem{theorem}{Theorem}
\newtheorem{corollary}[theorem]{Corollary}
\newcommand{\gt}{\gamma_t}

\title{Counterexamples to Graffiti.pc Conjecture 256\\
and a sharp total domination bound}
\author{}
\date{19 November 2026}

\begin{document}
\maketitle

\begin{abstract}
Graffiti.pc Conjecture 256 relates the total domination number of a
connected graph to the open neighborhood of its minimum-degree vertices
and the largest open neighborhood of an edge. A star with triangles
attached at its vertices gives an infinite family of counterexamples.
The same family shows that the conjectured coefficient \(2\) has optimal
replacement \(1\). A companion Lean~4 file verifies a 17-vertex
counterexample to the universal conjecture.
\end{abstract}

\section{Statement}

Let \(G=(V,E)\) be a finite simple connected graph with \(|V|\ge2\).
Write \(N(v)\) for the open neighborhood of \(v\), put
\(N(S)=\bigcup_{v\in S}N(v)\), and let \(A\) be the set of vertices of
minimum degree. Define
\[
M(G)=\max_{uv\in E}|N(u)\cup N(v)|.
\]
The number \(\gt(G)\) is the least size of a set \(D\subseteq V\)
such that \(N(D)=V\). In particular, vertices in \(D\) also require
a neighbor in \(D\).
Conjecture 256 in DeLaVina's \emph{Written on the Wall II}~\cite[pp.~148--149]{wow}
asserts that
\begin{equation}\label{eq:conjecture}
\gt(G)\ge \frac{2|N(A)|}{M(G)}.
\end{equation}
Under the catalogue's neighborhood convention, a member of \(A\)
belongs to \(N(A)\) whenever it has a neighbor in \(A\).

\section{Counterexamples}

For integers \(k\ge2\) and \(r\ge1\), start with a star with center \(c\)
and leaves \(x_1,\dots,x_k\). Attach one triangle at \(c\) and \(r\)
triangles at each \(x_i\). To attach a triangle at \(v\), add two fresh
vertices \(a,b\) and edges \(va,vb,ab\). Call the resulting graph \(G_{k,r}\).

\begin{theorem}\label{thm:family}
For \(G=G_{k,r}\),
\[
\gt(G)=k+1,\qquad |N(A)|=k(2r+1)+3,\qquad M(G)=2r+k+3.
\]
Consequently, \(G_{k,r}\) violates \eqref{eq:conjecture} exactly when
\(2r>k+3\).
\end{theorem}

\begin{proof}
The original star has no isolated vertices, and each new vertex is
adjacent to its attachment vertex. Thus the original star vertices
together totally dominate \(G\). Partition \(V\) into \(k+1\) groups
by putting each original vertex together with all the new vertices
attached to it. Every group contains a new vertex whose two neighbors
lie in that group. Therefore every total dominating set meets every
group, proving \(\gt(G)=k+1\).

The original center has degree \(k+2\), each original leaf has degree
\(2r+1\), and every new vertex has degree two. Thus \(A\) consists
exactly of the new vertices. Each original vertex has a new neighbor,
and each new vertex is adjacent to its triangle partner. Hence
\(N(A)=V\), whose size is \(k+1+2+2kr=k(2r+1)+3\).

On an original edge \(cx_i\), the neighborhoods of the two endpoints
are disjoint, so their union has size \((k+2)+(2r+1)\).
On an edge of an attached triangle the endpoints have exactly one
common neighbor, the third vertex of the triangle. The possible union
sizes for those edges are \(k+3\), \(2r+2\), and \(3\), all smaller
than \(2r+k+3\). This gives the claimed value of \(M(G)\).
Finally,
\[
2|N(A)|-\gt(G)M(G)=(k-1)(2r-k-3).
\]
Since \(k\ge2\) and \(M(G)>0\), the sign gives the stated criterion.
\end{proof}

\section{The optimal coefficient}

\begin{corollary}
The largest real \(\lambda\) for which
\(\gt(G)\ge \lambda|N(A)|/M(G)\) holds for every graph in the conjecture
is \(1\).
\end{corollary}

\begin{proof}
For every vertex \(u\), choose a neighbor \(v\). Then
\(|N(u)|\le |N(u)\cup N(v)|\le M(G)\). The neighborhoods of a minimum
total dominating set cover \(V\), so
\[
|N(A)|\le |V|\le \gt(G)M(G).
\]
Thus \(\lambda=1\) is valid.
To prove optimality, take \(r=k^2\) in Theorem~\ref{thm:family}.
The resulting ratios satisfy
\[
\frac{\gt(G_{k,k^2})M(G_{k,k^2})}{|N(A)|}
=\frac{(k+1)(2k^2+k+3)}{k(2k^2+1)+3}
\longrightarrow 1
\qquad(k\to\infty).
\]
Any \(\lambda>1\) therefore fails for sufficiently large \(k\).
\end{proof}

\section{Lean verification}

The companion file \texttt{Graffiti256.lean} gives the theorem
\texttt{Graffiti256.conjecture256\_false}. It encodes \(G_{2,3}\) on
\(\mathrm{Fin}\,17\), with original vertices numbered \(0,1,2\) and
center \(1\). Lean verifies the graph is simple and connected, verifies
that \(\{0,1,2\}\) totally dominates it, and checks
\(|N(A)|=17\) and \(M(G)=11\). The conjectured bound would imply
\(34\le 3\cdot11=33\) for this set, since any total dominating set
has size at least \(\gt(G)\). The formal result is the disproof of the
universal statement. The family and the optimal coefficient are
established by the mathematical arguments above.

\paragraph{Provenance.}
The catalogue snapshot updated on 5 November 2026 labels entry 256
open. The companion research notes record the search for earlier
resolutions and its limitations. AI assistance was used to develop
the argument, draft the manuscript, and write and check the Lean file.

\begin{thebibliography}{1}
\bibitem{wow}
E.~DeLaVina, \emph{Written on the Wall II: Conjectures of Graffiti.pc},
University of Houston--Downtown, catalogue snapshot printed
6 November 2026, pp.~148--149.
\href{https://www.uhd.edu/academics/sciences/mathematics-statistics/wowii/files/wowiiaii.pdf#page=149}{Catalogue, Conjecture 256}.
\end{thebibliography}

\end{document}
